\begin{summarybox}
	\textbf{\textcolor{CSummary}{一句话核心}}

	\textbf{椭圆上点到直线的距离最值，本质上由线性表达式 \(Ax+By\) 在椭圆上的值域决定。核心辅助量是}
	\[
		R=\sqrt{a^2A^2+b^2B^2}.
	\]

	\medskip

	\textbf{\textcolor{CSummary}{使用条件}}

	\begin{itemize}[leftmargin=*, itemsep=0.5em, topsep=0.4em]
		\item \textbf{条件 1（椭圆标准式）：}
		      椭圆方程为
		      \[
			      \frac{x^2}{a^2}+\frac{y^2}{b^2}=1,
		      \]
		      且
		      \[
			      a>b>0.
		      \]

		\item \textbf{条件 2（直线一般式）：}
		      直线方程为
		      \[
			      Ax+By+C=0,
		      \]
		      且 \(A,B\) 不全为零。

		\item \textbf{条件 3（辅助量）：}
		      \[
			      R=\sqrt{a^2A^2+b^2B^2}.
		      \]
	\end{itemize}

	\medskip

	\textbf{\textcolor{CSummary}{核心公式}}

	\[
		d_{\min}
		=
		\max\left\{0,\frac{|C|-R}{\sqrt{A^2+B^2}}\right\},
		\qquad
		d_{\max}
		=
		\frac{|C|+R}{\sqrt{A^2+B^2}}.
	\]

	\medskip

	\textbf{\textcolor{CSummary}{两种视角速记}}

	\begin{itemize}[leftmargin=*, itemsep=0.35em, topsep=0.3em]
		\item \textbf{椭圆动点 \(\to\) 固定直线：}
		      设 \((x,y)=(a\cos\theta,b\sin\theta)\)，
		      结合点到直线距离与辅助角公式求最值。

		\item \textbf{固定直线 \(\to\) 椭圆两侧：}
		      令切线法向量与 \((A,B)\) 成比例，得到两支撑切线
		      \[
			      Ax+By=R,\qquad Ax+By=-R.
		      \]
		      原直线到这两条切线的距离为
		      \[
			      D_+=\frac{|C+R|}{\sqrt{A^2+B^2}},\qquad
			      D_-=\frac{|C-R|}{\sqrt{A^2+B^2}}.
		      \]
		      最大值恒取 \(\max\{D_+,D_-\}\)；
		      最小值仅在 \(|C|>R\) 时取 \(\min\{D_+,D_-\}\)，
		      否则取 \(0\)。
	\end{itemize}

	\medskip

	\textbf{\textcolor{CSummary}{位置关系}}

	\[
		\left.
		\begin{aligned}
			|C|>R & \Longleftrightarrow \text{直线与椭圆相离},     \\
			|C|=R & \Longleftrightarrow \text{直线与椭圆相切},     \\
			|C|<R & \Longleftrightarrow \text{直线与椭圆相交于两个点}.
		\end{aligned}
		\right.
	\]

	\medskip

	\textbf{\textcolor{CSummary}{关键提醒}}

	\begin{itemize}[leftmargin=*, itemsep=0.5em, topsep=0.4em]
		\item \textbf{易错点（\(R\) 平方项）：}
		      计算 \(R\) 时务必注意
		      \[
			      a^2A^2+b^2B^2,
		      \]
		      也可以写成
		      \[
			      (aA)^2+(bB)^2.
		      \]

		\item \textbf{检查项（最小值取零）：}
		      当 \(|C|\le R\) 时，直线与椭圆有公共点，所以
		      \[
			      d_{\min}=0,
		      \]
		      不能把负数公式值当成最小距离。

		\item \textbf{取等点（不要误写成标准顶点）：}
		      最大值或最小值一般在法向支撑点处取得，而不一定在椭圆的四个标准顶点处取得。对应支撑点为
		      \[
			      P_+=\left(\frac{a^2A}{R},\frac{b^2B}{R}\right),
			      \qquad
			      P_-=-P_+.
		      \]

		\item \textbf{几何意义（区分 \(R\) 与真实距离）：}
		      \(R\) 是 \(Ax+By\) 在椭圆上的最大绝对值；
		      \[
			      \frac{R}{\sqrt{A^2+B^2}}
		      \]
		      才是椭圆沿直线法向方向的真实投影半宽。
	\end{itemize}

\end{summarybox}